\documentclass{article}
\usepackage{amsmath}
\usepackage{amssymb}

\title{D-S1-1: Period of the simple pendulum to order $\theta_0^2$}
\date{}

\begin{document}
\maketitle

\section*{Target and conventions}

Simple pendulum: point mass $m$ on a massless rigid rod of length $L$, no
friction, uniform gravity $g$. Following \texttt{conventions.md}, $\theta(t)$
is the angle of the rod from the downward vertical, the equation of motion is
\begin{equation}
  \ddot\theta + \omega_0^2 \sin\theta = 0, \qquad \omega_0 = \sqrt{g/L},
  \qquad T_0 = \frac{2\pi}{\omega_0} = 2\pi\sqrt{L/g},
  \label{eq:eom}
\end{equation}
$\theta_0 \in (0,\pi)$ is the amplitude (maximum angle reached), and the
complete elliptic integral of the first kind is written with the
\emph{modulus} $k$:
\begin{equation}
  K(k) = \int_0^{\pi/2} \frac{d\phi}{\sqrt{1 - k^2 \sin^2\phi}} .
  \label{eq:Kdef}
\end{equation}
Target: the coefficient $c_2$ in
$T(\theta_0) = T_0\,(1 + c_2\,\theta_0^2 + O(\theta_0^4))$.

\section{Exact period}

\paragraph{Energy integral.}
Multiplying \eqref{eq:eom} by $\dot\theta$ and integrating once in time,
with $\dot\theta = 0$ at $\theta = \theta_0$ (a turning point),
\begin{equation}
  \tfrac12 \dot\theta^2 - \omega_0^2 \cos\theta = - \omega_0^2 \cos\theta_0
  \quad\Longrightarrow\quad
  \dot\theta^2 = 2\omega_0^2 \,(\cos\theta - \cos\theta_0).
  \label{eq:energy}
\end{equation}
(This is conservation of $\tfrac12 m L^2\dot\theta^2 - m g L\cos\theta$,
divided by $mL^2$; the mass drops out.) For $0<\theta_0<\pi$ the motion is a
libration between $-\theta_0$ and $\theta_0$, and these are the only turning
points, since $\dot\theta = 0$ requires $\cos\theta = \cos\theta_0$.

\paragraph{Quarter period.}
The equation \eqref{eq:eom} is invariant under $\theta \to -\theta$ and under
time reversal $t \to -t$. Hence the four legs $\theta_0 \to 0 \to -\theta_0
\to 0 \to \theta_0$ take equal times, and on the leg from $0$ to $\theta_0$,
$\dot\theta > 0$. With \eqref{eq:energy},
\begin{equation}
  T(\theta_0) = 4\int_0^{\theta_0} \frac{d\theta}{\dot\theta}
  = \frac{4}{\omega_0} \int_0^{\theta_0}
    \frac{d\theta}{\sqrt{2(\cos\theta - \cos\theta_0)}} .
  \label{eq:Tint}
\end{equation}
The integrand has an integrable singularity at $\theta = \theta_0$: there
$\cos\theta - \cos\theta_0 \simeq \sin\theta_0\,(\theta_0 - \theta)$ with
$\sin\theta_0 > 0$, so the integrand behaves as
$(\theta_0-\theta)^{-1/2}$.

\paragraph{Half-angle form.}
Using $\cos\theta = 1 - 2\sin^2(\theta/2)$ and writing
\begin{equation}
  k \equiv \sin(\theta_0/2) \in (0,1),
  \label{eq:k}
\end{equation}
one has $\cos\theta - \cos\theta_0 = 2\,\bigl(k^2 - \sin^2(\theta/2)\bigr)$,
so \eqref{eq:Tint} becomes
\begin{equation}
  T(\theta_0) = \frac{2}{\omega_0}\int_0^{\theta_0}
    \frac{d\theta}{\sqrt{k^2 - \sin^2(\theta/2)}} .
  \label{eq:Thalf}
\end{equation}

\paragraph{Substitution.}
Set $\sin(\theta/2) = k\sin\phi$ with $\phi \in [0,\pi/2]$. As $\theta$ runs
from $0$ to $\theta_0$, $\sin(\theta/2)$ runs monotonically from $0$ to $k$
(because $\theta/2 \in [0,\pi/2)$), so $\phi$ runs from $0$ to $\pi/2$.
Differentiating,
\begin{equation}
  \tfrac12 \cos(\theta/2)\, d\theta = k\cos\phi \, d\phi
  \quad\Longrightarrow\quad
  d\theta = \frac{2k\cos\phi\, d\phi}{\sqrt{1 - k^2\sin^2\phi}},
  \label{eq:dtheta}
\end{equation}
where $\cos(\theta/2) = \sqrt{1-\sin^2(\theta/2)} = \sqrt{1-k^2\sin^2\phi}$
(positive since $\theta/2\in[0,\pi/2)$). Also
$\sqrt{k^2 - \sin^2(\theta/2)} = \sqrt{k^2 - k^2\sin^2\phi} = k\cos\phi$.
Substituting into \eqref{eq:Thalf}, the factors $k\cos\phi$ cancel:
\begin{equation}
  T(\theta_0) = \frac{2}{\omega_0}\int_0^{\pi/2}
    \frac{2\,d\phi}{\sqrt{1-k^2\sin^2\phi}}
  = \frac{4}{\omega_0}\, K(k).
  \label{eq:Texact0}
\end{equation}
With $T_0 = 2\pi/\omega_0$, the exact period is
\begin{equation}
  \boxed{\;
  T(\theta_0) = \frac{2\,T_0}{\pi}\, K\!\bigl(\sin(\theta_0/2)\bigr)
  = 4\sqrt{\frac{L}{g}}\; K\!\bigl(\sin(\theta_0/2)\bigr)
  \;}
  \label{eq:Texact}
\end{equation}
with $K$ in the modulus convention \eqref{eq:Kdef}. (In the parameter
convention $m = k^2$, the argument is $\sin^2(\theta_0/2)$.)

\section{Expansion in the amplitude}

\paragraph{Series for $K$.}
For $|x| < 1$, the binomial series gives
\begin{equation}
  (1-x)^{-1/2} = \sum_{n=0}^{\infty} \binom{2n}{n}\frac{x^n}{4^n}.
  \label{eq:binom}
\end{equation}
With $x = k^2\sin^2\phi$ and $0 \le x \le k^2 < 1$ on
$\phi \in [0,\pi/2]$, the series converges uniformly in $\phi$ (Weierstrass
test with majorant $\sum \binom{2n}{n} k^{2n}/4^n < \infty$), so it may be
integrated term by term. Using the Wallis integral
\begin{equation}
  \int_0^{\pi/2} \sin^{2n}\phi\, d\phi = \frac{\pi}{2}\,
  \frac{(2n-1)!!}{(2n)!!} = \frac{\pi}{2}\binom{2n}{n}\frac{1}{4^n},
  \label{eq:wallis}
\end{equation}
one gets
\begin{equation}
  K(k) = \frac{\pi}{2}\sum_{n=0}^{\infty}
    \left[\binom{2n}{n}\frac{1}{4^n}\right]^2 k^{2n}
  = \frac{\pi}{2}\left(1 + \frac{1}{4}k^2 + \frac{9}{64}k^4
    + \frac{25}{256}k^6 + O(k^8)\right).
  \label{eq:Kseries}
\end{equation}
The coefficients are $(\tfrac{2}{4})^2 = \tfrac14$,
$(\tfrac{6}{16})^2 = \tfrac{9}{64}$, $(\tfrac{20}{64})^2 = \tfrac{25}{256}$.
Inserting \eqref{eq:Kseries} into \eqref{eq:Texact},
\begin{equation}
  \frac{T(\theta_0)}{T_0} = 1 + \frac{1}{4}k^2 + \frac{9}{64}k^4
    + \frac{25}{256}k^6 + O(k^8),
  \qquad k^2 = \sin^2(\theta_0/2).
  \label{eq:Tk}
\end{equation}

\paragraph{Series for $k^2$.}
From $\sin^2 x = (1-\cos 2x)/2 = x^2 - \tfrac{1}{3}x^4 + \tfrac{2}{45}x^6
+ O(x^8)$ with $x = \theta_0/2$,
\begin{equation}
  k^2 = \sin^2(\theta_0/2) = \frac{\theta_0^2}{4} - \frac{\theta_0^4}{48}
    + \frac{\theta_0^6}{1440} + O(\theta_0^8).
  \label{eq:k2}
\end{equation}
Hence $k^4 = \theta_0^4/16 - \theta_0^6/96 + O(\theta_0^8)$ and
$k^6 = \theta_0^6/64 + O(\theta_0^8)$.

\paragraph{Coefficients.}
Collecting powers of $\theta_0$ in \eqref{eq:Tk}:
\begin{align}
  \theta_0^2:&\quad \frac14\cdot\frac14 = \frac{1}{16}, \label{eq:c2}\\
  \theta_0^4:&\quad \frac14\cdot\Bigl(-\frac{1}{48}\Bigr)
     + \frac{9}{64}\cdot\frac{1}{16}
     = -\frac{1}{192} + \frac{9}{1024} = \frac{-16 + 27}{3072}
     = \frac{11}{3072}, \label{eq:c4}\\
  \theta_0^6:&\quad \frac14\cdot\frac{1}{1440}
     + \frac{9}{64}\cdot\Bigl(-\frac{1}{96}\Bigr)
     + \frac{25}{256}\cdot\frac{1}{64}
     = \frac{128 - 1080 + 1125}{737280} = \frac{173}{737280}. \label{eq:c6}
\end{align}
Therefore
\begin{equation}
  \boxed{\;
  T(\theta_0) = T_0\left(1 + \frac{1}{16}\,\theta_0^2
    + \frac{11}{3072}\,\theta_0^4 + O(\theta_0^6)\right),
  \qquad c_2 = \frac{1}{16}
  \;}
  \label{eq:result}
\end{equation}
(and $c_6 = 173/737280$ from \eqref{eq:c6}, used only in the sanity check
below). The expansion in $\theta_0$ is a convergent power series near
$\theta_0 = 0$: $K$ is analytic in $k^2$ for $|k^2|<1$, and
$k^2 = \sin^2(\theta_0/2)$ is an entire function of $\theta_0$ vanishing at
$0$, so the composition is analytic in a neighbourhood of $\theta_0 = 0$.

\section{Independent check of $c_2$: Poincar\'e--Lindstedt}

This check does not use the exact period. Expand
$\sin\theta = \theta - \theta^3/6 + O(\theta^5)$ and introduce
$\tau = \omega t$, where $\omega = 2\pi/T$ is the (unknown) angular frequency,
so that $\theta$ is $2\pi$-periodic in $\tau$. Write
$\theta = A\,u(\tau)$, $u = u_1 + A^2 u_3 + O(A^4)$,
$\omega^2 = \omega_0^2(1 + A^2 w + O(A^4))$, with initial conditions
$u(0) = 1$, $u'(0) = 0$. Equation \eqref{eq:eom} becomes
\begin{equation}
  (1 + A^2 w)\,(u_1'' + A^2 u_3'') + u_1 + A^2 u_3 - \frac{A^2}{6}u_1^3
  = O(A^4).
  \label{eq:PL}
\end{equation}
Order $A^0$: $u_1'' + u_1 = 0$, $u_1(0)=1$, $u_1'(0)=0$, so
$u_1 = \cos\tau$. Order $A^2$, with $u_1'' = -\cos\tau$ and
$\cos^3\tau = \tfrac34\cos\tau + \tfrac14\cos 3\tau$:
\begin{equation}
  u_3'' + u_3 = w\cos\tau + \frac{1}{6}\cos^3\tau
  = \Bigl(w + \frac18\Bigr)\cos\tau + \frac{1}{24}\cos 3\tau .
  \label{eq:PL3}
\end{equation}
Periodicity of $u_3$ requires the resonant ($\cos\tau$) term to vanish:
$w = -1/8$. Then $\omega = \omega_0\,(1 - A^2/16 + O(A^4))$ and
\begin{equation}
  T = \frac{2\pi}{\omega} = T_0\left(1 + \frac{A^2}{16} + O(A^4)\right).
  \label{eq:PLresult}
\end{equation}
The solution of \eqref{eq:PL3} with $u_3(0) = u_3'(0) = 0$ is
$u_3 = (\cos\tau - \cos 3\tau)/192$, so $\theta(0) = A$ and
$\dot\theta(0) = 0$: $\tau = 0$ is a turning point. By the remark after
\eqref{eq:energy}, every turning point has $|\theta| = \theta_0$, so
$A = \theta_0$. Hence \eqref{eq:PLresult} gives $c_2 = 1/16$, in agreement
with \eqref{eq:c2}.

\section{Sanity checks}

\begin{itemize}
  \item \emph{Small amplitude.} As $\theta_0 \to 0$, $k \to 0$ and
    $K(0) = \pi/2$, so \eqref{eq:Texact} gives $T \to T_0$.
  \item \emph{Large amplitude.} As $\theta_0 \to \pi$, $k \to 1$ and $K(k)$
    diverges logarithmically, so $T \to \infty$: the pendulum takes an
    infinite time to approach the unstable upright equilibrium.
  \item \emph{Sign.} $c_2 > 0$. Since $\sin\theta < \theta$ for
    $0<\theta<\pi$, the restoring torque is weaker than in the linearised
    problem, so the period grows with amplitude.
  \item \emph{Parity.} \eqref{eq:eom} is invariant under
    $\theta\to-\theta$, so $T$ is even in $\theta_0$; only even powers
    appear in \eqref{eq:result}.
  \item \emph{Dimensions.} $K$ and $\theta_0$ (radians) are dimensionless,
    $\sqrt{L/g}$ is a time; $c_2$ is a pure number. The mass does not enter.
  \item \emph{Two methods.} $c_2 = 1/16$ is obtained both from the exact
    period (\S2) and from Poincar\'e--Lindstedt (\S3), which share only the
    equation of motion.
  \item \emph{Hand evaluation at $\theta_0 = 0.1$.} $k = \sin(0.05) =
    0.0499791693$, $k^2 = 0.0024979174$. Summing \eqref{eq:Tk} through
    $k^8$ gives $T/T_0 - 1 = 6.2535831\times10^{-4}$, i.e.
    $(T/T_0 - 1)/\theta_0^2 = 0.06253583$. The truncated series
    $1/16 + \tfrac{11}{3072}\theta_0^2 + \tfrac{173}{737280}\theta_0^4 =
    0.06253583$ agrees to the eight digits kept, and the deviation from
    $1/16 = 0.0625$ is accounted for by the $\theta_0^4$ term. This was
    done by hand (no code was run).
\end{itemize}

\end{document}
